
%=========================================================
\section{Pricing-Consistency Diagnostics}
\label{sec:pricing_diagnostics}
%=========================================================

The Monte Carlo pricing pipeline defined above relies on two identities being satisfied (up to MC error) by the trained model. First, the model-implied discount curve must be reproduced by pathwise discounting,
\begin{equation}
  \mathbb{E}^{\mathbb Q}\!\bigl[D(0,T)\bigr]
  \;=\;
  P(z_0, T),
  \qquad T>0,
  \label{eq:zero-martingale}
\end{equation}
and, more generally, the discounted-bond martingale identity
\begin{equation}
  \mathbb{E}^{\mathbb Q}\!\bigl[D(0,t)\,P(z_t, T-t)\bigr]
  \;=\;
  P(z_0, T),
  \qquad 0\le t\le T,
  \label{eq:bond-martingale}
\end{equation}
must hold for every intermediate time $t$. Second, the forward swap rate computed from the time-zero curve must equal the annuity-numeraire expectation of the simulated rate at expiry, as already stated in \eqref{eq:swap-martingale}, or equivalently in risk-neutral form
\begin{equation}
  F(0)
  \;=\;
  \frac{\mathbb{E}^{\mathbb Q}\!\left[D(0,T_\mu)\,\mathcal A(T_\mu)\,F(T_\mu)\right]}
       {\mathbb{E}^{\mathbb Q}\!\left[D(0,T_\mu)\,\mathcal A(T_\mu)\right]}.
  \label{eq:swap-martingale-Q}
\end{equation}
A violation of \eqref{eq:zero-martingale}--\eqref{eq:swap-martingale-Q} implies that the ATM strike, the annuity, or the pathwise discount used by the Monte Carlo estimator are mutually inconsistent, so any Bachelier implied volatility extracted by inverting \eqref{eq:atm_bachelier_price} will be biased independently of the diffusion calibration.

We assess all three identities for the stable model at every latent dimension $d\in\{2,3,4\}$ in the EUR-only setting, using \(N_{\mathrm{MC}}=10{,}000\) antithetic paths simulated to a ten-year horizon. To separate causes from symptoms, the simulation-based diagnostics in Section~\ref{subsec:simulation_diagnostics_results} are paired with a purely static analysis of the trained checkpoint (Section~\ref{subsec:static_diagnostics_results}), which evaluates the encoder, decoder, and short-rate functions on the encoded training cloud only.

%---------------------------------------------------------
\subsection{Discounted-bond martingale and zero-coupon repricing}
\label{subsec:simulation_diagnostics_results}
%---------------------------------------------------------

Table~\ref{tab:zero_repricing} reports the yield-error of \eqref{eq:zero-martingale}, expressed in basis points, for maturities $T\in\{1,2,3,5,7,10\}$ years and the three latent dimensions. None of the checkpoints satisfies the identity beyond $T=3$\,Y. At $T=10$\,Y the absolute yield error grows to $139$, $1\,388$, and $396$ bps for dim $2$, $3$, and $4$ respectively; the Monte Carlo discount in dim $3$ and dim $4$ even exceeds one, which is not a valid discount factor. Doubling the path count from $4\,000$ to $10\,000$ moved the dim 4 error by less than $3$ bps, confirming that the gap is structural model bias rather than sampling noise.

\begin{table}[ht]
  \centering
  \caption{Zero-coupon repricing error
    \(\hat y_{\mathrm{MC}}(T)-y_{\mathrm{dec}}(T)\) in basis points
    of yield, EUR curve, \(N_{\mathrm{MC}}=10{,}000\) paths,
    monthly Euler.}
  \label{tab:zero_repricing}
  \begin{tabular}{rrrr}
    \toprule
    $T$ (yr) & dim 2 & dim 3 & dim 4 \\
    \midrule
       1 & $+123.4$  &  $+749.5$ & $+29.3$  \\
       2 & $+72.9$   &  $+601.5$ & $-23.9$  \\
       3 & $+32.4$   &  $+393.0$ & $-85.4$  \\
       5 & $-34.3$   &  $-167.6$ & $-193.5$ \\
       7 & $-85.6$   &  $-706.8$ & $-283.5$ \\
      10 & $-138.7$  & $-1387.7$ & $-395.7$ \\
    \bottomrule
  \end{tabular}
\end{table}

The full discounted-bond martingale \eqref{eq:bond-martingale} fails in the same way at every intermediate time. For example, in dim 4 the estimator $\hat M(t,10)=\mathbb E[D(0,t)P(z_t,10-t)]$ is $0.7647$ at $t=0.25$\,Y, falls to $0.6733$ at $t=2$\,Y, and rises to $1.0534$ at $t=10$\,Y, instead of remaining flat at the time-zero decoder value $P(z_0,10)=0.7092$. A truly risk-neutral trained model would produce a flat trajectory up to MC noise.

The forward-swap identity \eqref{eq:swap-martingale} is similarly violated. For the $5\text{Y}\times 5\text{Y}$ structure the decoder forward $F(0)$ is $434.7$, $440.6$, and $434.7$ bps in dim $2$, $3$, and $4$, while the annuity-weighted Monte Carlo estimate of $\mathbb E^{\mathcal A}[F(T_\mu)]$ is $324.5$, $3\,069.5$, and $247.0$ bps respectively, even after discarding paths whose discount curve violates the bounds \(0<P\le1\). Through \eqref{eq:bias_parity} this translates into a forward-centring bias $\widehat b$ of roughly $-110$, $+2\,629$, and $-188$ bps; the mis-centring of the strike therefore biases every Bachelier implied volatility extracted from \eqref{eq:straddle_normal_vol}, irrespective of how well the diffusion is calibrated.

The decoder itself produces an invalid discount curve on a substantial fraction of simulated paths. By construction the training-point curves satisfy the no-arbitrage PDE up to machine precision, but the simulated cloud exits the training support within months. At $t=10$\,Y the share of paths with \(P(z_t,\tau)\le1\) and monotone in $\tau$ is $67\%$, $0.1\%$, and $19\%$ in dim $2$, $3$, and $4$; the share with at least one entry exceeding the double-precision range or otherwise non-finite is $6\%$, $42\%$, and $46\%$ in the same three dimensions.

\paragraph{Closed-form residual via the Riccati ODE.}
To evaluate the no-arbitrage PDE residual without relying on the
decoder's own output, we exploit the affine structure of the model.
Writing the decoder factor as $G(z,\tau)$, the bond price admits the
representation \citep{Poulsen2024}
\[
  P(z,\tau)\;=\;\exp\!\bigl(-A(\tau)-B(\tau)^{\!\top}z\bigr),
\]
where the coefficient functions $A:\mathbb{R}_{+}\!\to\!\mathbb{R}$
and $B:\mathbb{R}_{+}\!\to\!\mathbb{R}^{d}$ solve the coupled Riccati
system
\begin{equation}
  \frac{dB}{d\tau} \;=\; \alpha(\tau)\,B(\tau) + \beta(\tau),
  \qquad
  \frac{dA}{d\tau} \;=\; \gamma(\tau)\,B(\tau)^{2},
  \qquad A(0)=0,\;B(0)=0,
  \label{eq:ab_riccati}
\end{equation}
with coefficients
\[
  \alpha \;=\; \frac{-\partial_\tau G + \nabla_{\!z}G^{\!\top}\mu(z) + \tfrac12\mathrm{Tr}\bigl[\sigma^{\!\top}\nabla_{\!z}^{2}G\,\sigma\bigr]}{G},
  \quad
  \beta  \;=\; \frac{\tilde r(z)}{G},
  \quad
  \gamma \;=\; \tfrac12\bigl\lVert\sigma(z)^{\!\top}\nabla_{\!z}G\bigr\rVert^{2}.
\]
The derivatives $\partial_\tau G$, $\nabla_{\!z}G$ and
$\mathrm{Tr}[\sigma^{\!\top}\nabla_{\!z}^{2}G\,\sigma]$ are obtained
analytically by forward-mode automatic differentiation through the
decoder network, vectorised over the maturity grid
(\texttt{Code/utils/ode.py}). System~\eqref{eq:ab_riccati} is then
integrated by the four-stage \(\tfrac{3}{8}\)-rule Runge--Kutta scheme
on the same $\tau$-grid as the decoder. The PDE residual reported
below is
\[
  R_\tau(z)
  \;=\;
  \partial_\tau P(z,\tau) + \tilde r(z)\,P(z,\tau)
  -\;\nabla_{\!z}P^{\!\top}\mu(z)
  -\;\tfrac12\mathrm{Tr}\bigl[\sigma^{\!\top}\nabla_{\!z}^{2}P\,\sigma\bigr],
\]
evaluated by comparing the decoder's $P(z,\tau)$ against the
exponential-affine value implied by the integrated $A(\tau),B(\tau)$
trajectory. Because the integration runs entirely in latent space and
uses no Monte Carlo paths, $R_\tau$ isolates pure decoder
extrapolation error.

The Poulsen PDE residual $\max_\tau\lvert R_\tau\rvert$ rises from
$\sim 10^{-17}$ at training points to $\mathcal O(1)$ on the $99$-th
percentile of $z_t$ in dim $3$ and dim $4$. This is the proximate
cause of the martingale failure: the decoder is no longer
arbitrage-free at the points where the Monte Carlo estimator
evaluates it.

The simulation-based source of the failure is identified by the
equilibrium of the linear drift $\mu(z)=Mz+N$. Setting
\(\mu(z^\ast)=0\) yields \(z^\ast=-M^{-1}N\), which is
\((-0.19,-0.06)\) for dim $2$, \((-242.6,-151.7,+79.1)\) for dim $3$,
and \((-8.5,+3.8,+1.7,-0.6)\) for dim $4$. The training-cloud means
are \((-0.13,+0.02)\), \((+0.03,-0.05,+0.05)\), and
\((+0.07,-0.00,+0.03,-0.01)\) respectively. Only the two-factor
model places the stationary point inside the encoded support.
Furthermore, the slowest mean-reversion half-life
$-\log 2/\Re(\lambda_{\min})$ is $0.52$ yr in dim $2$, but
$684$ yr in dim $3$ and $423$ yr in dim $4$, because the
parameterisation \(M=-(V^{\!\top}V+\varepsilon I)\) admits a
near-zero eigenvalue when one row of $V$ collapses during training.
That direction behaves as a Brownian motion under $\mathbb Q$ over
the ten-year horizon, driving $z_t$ away from the data manifold.

%---------------------------------------------------------
\subsection{Static consistency of the trained checkpoint}
\label{subsec:static_diagnostics_results}
%---------------------------------------------------------

To verify that the failures above reflect the model itself and not
the Euler--Maruyama scheme, we evaluate the same checkpoints on a
purely static set of points: the $170$ encoded EUR training curves
and the $70$ held-out EUR curves used for out-of-sample analysis in
Section~\ref{sec:oos_results}.

The autoencoder reconstructs the training swap rates accurately in
all three dimensions, with global RMSE $8.5$, $11.4$, and $6.0$ bps
for dim $2$, $3$, and $4$. The out-of-sample RMSE is $11.6$, $10.2$,
and $8.1$ bps, so higher latent dimension does improve held-out
reconstruction. The Poulsen PDE residual at training points is at
the level of machine precision in every dimension, and the
swap-reduction identity
\(F=(1-P_T)/\sum_k P_k\) holds to zero. The model is therefore
internally consistent wherever it has been fit.

The dimensional improvement of the reconstruction is, however,
\emph{not} reflected in the latent geometry. The empirical
explained-variance ratios of the encoded training cloud are
\((0.88,\,0.12)\) in dim $2$,
\((0.95,\,0.05,\,3\!\times\!10^{-4})\) in dim $3$, and
\((0.86,\,0.13,\,0.016,\,9\!\times\!10^{-7})\) in dim $4$. The
fourth axis of the dim-4 model carries no variance, and the third
axis of dim $3$ carries three orders of magnitude less than the
second. On the held-out set the standard deviations of those
collapsed axes are even smaller. The encoder thus fits the data
using only two or three latent directions in practice, while
the remaining dimensions act as additional neural-network parameters
that widen the encoder and decoder linear maps.

This decoupling explains why higher latent dimension simultaneously
improves reconstruction and degrades simulation. The trained drift
matrix $M$ inherits the encoder's rank deficiency through the
parameterisation \(M=-(V^{\!\top}V+\varepsilon I)\): when one row of
$V$ collapses, the corresponding eigenvalue of $M$ is bounded only by
the regularisation floor $\varepsilon=10^{-3}$, producing the
near-zero eigenvalues reported in
Section~\ref{subsec:simulation_diagnostics_results}.

The reverse direction is verified by a static perturbation test.
Sampling \(z=\bar z+ k\,\sigma_z u\) with $u$ uniform on the unit
sphere and $\bar z,\sigma_z$ the per-axis mean and standard
deviation of the training cloud, we record the fraction of perturbed
states whose decoded curve is monotone with \(0<P\le1\).
Table~\ref{tab:static_perturbation} reports this share as a function
of the perturbation radius $k$.

\begin{table}[ht]
  \centering
  \caption{Share of static perturbations
    \(z=\bar z+k\sigma_z u\) for which the decoder produces a
    monotone discount curve with \(0<P\le1\). No SDE is simulated.}
  \label{tab:static_perturbation}
  \begin{tabular}{rrrr}
    \toprule
    $k$ ($\sigma_z$) & dim 2 & dim 3 & dim 4 \\
    \midrule
    0 & $100$\,\% & $100$\,\% & $100$\,\% \\
    1 & $72.3$\,\% & $30.4$\,\% & $71.2$\,\% \\
    2 & $50.7$\,\% & $21.3$\,\% & $55.6$\,\% \\
    3 & $36.9$\,\% & $18.0$\,\% & $47.6$\,\% \\
    \bottomrule
  \end{tabular}
\end{table}

Already at a one-standard-deviation perturbation, between $30\%$
and $70\%$ of decoded curves fail the validity check, before any
diffusion is introduced. The Monte Carlo failure of
Section~\ref{subsec:simulation_diagnostics_results} is therefore not
a consequence of the SDE; it is a consequence of the decoder being
fit only on a narrow data cloud and extrapolating outside it.

%---------------------------------------------------------
\subsection{Implications for the implied volatility surface}
\label{subsec:diagnostics_implications}
%---------------------------------------------------------

The combined evidence has direct consequences for the
model-market gap reported in Section~\ref{sec:market_comparison}.
Three pricing-pipeline inputs are inconsistent at once: the
time-zero forward swap rate $F(0)$ comes from the decoder, the
strike $K$ is set to that same value, while the payoff
$\mathcal A(T_\mu)(F(T_\mu)-K)^+$ is built from simulated curves
that no longer respect
\eqref{eq:zero-martingale}--\eqref{eq:swap-martingale-Q}. The
payer--receiver averaging in \eqref{eq:straddle_normal_vol} removes
the first-order directional bias \eqref{eq:bias_parity}, but two
effects remain. First, the $5\!\times\!5$ forward swap rate is
mis-centred by roughly $190$\,bps between the decoder and the
annuity-weighted Monte Carlo estimate
(Section~\ref{subsec:simulation_diagnostics_results}); that
mis-centring alone is enough to inflate a $70$-bp ATM
volatility by an order of magnitude through
\eqref{eq:atm_bachelier_price}. Second, the variance of the
simulated $F(T_\mu)$ across paths is dominated by the tail of paths
whose decoded discount curve violates $0<P\le1$; these contribute a
spurious volatility term that does not exist in the market.

These two effects act in the same direction and account
quantitatively for the over-pricing observed in
Section~\ref{sec:market_comparison}: the average market-vs-model
gap of $185$\,bps at $10\!\times\!10$ and $678$\,bps at
$1\!\times\!1$ is consistent with the $400$\,bp zero-coupon yield
error at $T=10$\,Y reported in Table~\ref{tab:zero_repricing}
amplified by the inverse-expiry factor $1/\sqrt{T_\mu}$ that appears
in the Bachelier inversion \eqref{eq:bachelier_inversion}. The shape
inversion, namely that the model volatility \emph{decreases} with
expiry while the market volatility \emph{increases}, is the
signature of the inverse-expiry amplification: short expiries divide
by a small $\sqrt{T_\mu}$ and so suffer the most from a fixed bias
in $\widehat V_0$. Both phenomena are therefore explained by the
simulation and static diagnostics, not by any defect in the
Bachelier inversion itself.

The dim-$2$ checkpoint is the least pathological in every category
above, even though it is dominated by dim $3$ and dim $4$ in OOS
reconstruction. The two failure modes therefore have to be
addressed by changes to the training objective rather than by
choosing a higher latent dimension. The static perturbation test
and the encoder PCA collapse together suggest two targeted
modifications: (i) a latent whitening penalty
$\sum_i\bigl(\log\sigma_{z_i}-\log\sigma_{z_0}\bigr)^{\!2}$ that
prevents axes from collapsing and so removes the near-zero
eigenvalue of $M$, and (ii) a perturbation-based no-arbitrage
penalty
$\mathbb E_{\varepsilon}\!\bigl[\max_\tau\lvert R_\tau(z+\varepsilon,\tau)\rvert^{2}\bigr]$
that forces the decoder to remain arbitrage-free in a neighbourhood
of the training cloud. Implementing either is left to future
training experiments and falls outside the scope of this thesis.
